MHC-I presentation can vary among tissues even when the presenting allele and source protein are fixed. Existing predictors largely estimate tissue-blind peptide--MHC compatibility and therefore do not directly address this tissue-associated preference. We formulate a matched-pair ranking problem: given two peptides from the same parent protein under the same presenting restriction, determine which has stronger presentation evidence in a specified tissue. The benchmark contains 157 human tissue--HLA combinations and 24 mouse tissue--H-2 combinations, where H-2 is the mouse major histocompatibility complex; each contains more than 200 matched peptide pairs. TissuePMHC combines a peptide representation learned across all human tasks with an allele-restricted representation, while a separately trained shared-learning model is used for mouse. On saved internal test sets, the equal-weight mean AUROC/AUPRC is 0.8448/0.8348 in human and 0.8562/0.8506 in mouse. When every held-out peptide identity is excluded from all training tasks, AUROC decreases to 0.7652 and 0.7529, respectively. Removing tissue identity produces a further substantial decrease, showing that general peptide--MHC compatibility alone does not reproduce the complete ranking signal. However, the full human and mouse architectures do not consistently outperform their strongest simpler controls after peptide separation. These results support learnability of tissue-conditioned peptide ranking for previously represented tissue--MHC combinations and quantify the optimism associated with cross-task peptide reuse. They do not establish a causal tissue-processing mechanism, generalization to unseen tissues or alleles, independent-cohort validity, or clinical utility.
